\chapter{Measuring a complete kernel}
\MARSChapterHeading{Measuring a complete kernel}
\label{revision:measurement}

An implementation can perform fewer memory misses while executing more
instructions. Those observations answer different questions. Before choosing a
version, establish that both compute the required result, specify where counting
begins and ends, and identify which costs the measurement represents.

The complete \texttt{matrix-measure.asm} program in \texttt{revision/labs/files}
provides an ordinary matrix product and a blocked product. Both are scalar MIPS64
programs implementing original project \texttt{022-blocked-matmul}. They use the
same arithmetic interface and select one method per launch. This chapter begins
the advanced sequence with functions, nested loops, full-width arithmetic and the
cache model from Chapter~\ref{revision:locality}; MSA knowledge is not required yet.

You will derive a product, identify the state carried between reduction blocks,
account for native memory notices, and interpret controlled measurements. The
independent exercise applies the same dot-product reasoning to a matrix-vector
interface. Complete source is supplied for observation, while short fragments
below expose the decisions that matter to the comparison.

\section{Establish the result before counting its cost}

For an $M\times K$ matrix $A$ and a $K\times N$ matrix $B$, their product has
$M$ rows and $N$ columns. Each output uses one row of $A$ and one column of $B$:
\[
  C[i,j]=\sum_{q=0}^{K-1} A[i,q]B[q,j].
\]
The index $q$ identifies corresponding entries to multiply and add; it does not
select another output cell. This operation is a \emph{dot product}. It differs
from multiplying entries at equal array indices or transposing either input.

The default source contains
\[
  A=\begin{bmatrix}1&2&3\\4&5&6\end{bmatrix},\qquad
  B=\begin{bmatrix}7&8\\9&10\\11&12\end{bmatrix}.
\]
The first result is $1\cdot7+2\cdot9+3\cdot11=58$. The next column uses
$8,10,12$, giving 64. The complete product is
\[
  C=\begin{bmatrix}58&64\\139&154\end{bmatrix}.
\]
Write these four sums before running the program. A returned count of four says
how many output cells the function reports; it does not prove their values.

Input elements are signed values from $-32768$ through 32767 stored in four-byte
words. Output elements are signed eight-byte integers. Each product has magnitude
at most $2^{30}$, and at most sixteen products contribute to a result, so its
magnitude is bounded by $2^{34}$. Doubleword arithmetic can represent every
allowed sum exactly. There is no fixed-point scaling, saturation or intended
wraparound in this contract.

Two products already expose the word boundary: with both input pairs equal to
$-32768$, the sum is $2^{31}=2147483648$. Its low word is hexadecimal
\texttt{80000000}, while its high word is zero. Interpreting only the low word as
a signed integer would report a negative value. The accumulator and output store
must preserve the complete result. The worked code uses \texttt{dmul} and
\texttt{daddu}; a signed word multiply can also produce these bounded individual
products correctly, but word accumulation can discard a required upper bit.

Both functions receive full pointers to $A,B,C$ in
\texttt{\$a0}/\texttt{\$a1}/\texttt{\$a2}. The pointer in \texttt{\$a3} names
four words $M,K,N,T$, where dimensions range from zero through sixteen and tile
edge $T$ from one through eight. Arrays are packed row major. The active input
windows have $MK$ and $KN$ words; output has $MN$ doublewords. Passed pointers
may identify interior, disjoint windows and may exceed 32 bits.

If $M=0$ or $N=0$, there is no output and no matrix access. If $K=0$ with a
nonempty output, each empty sum is zero: write all $MN$ zero doublewords without
input reads. Initial output bytes are arbitrary. Preserve inputs, descriptor,
guards and inactive output storage, and initialize an output byte before reading
it as a partial sum. Preserve the full saved registers, frame pointer and stack
pointer; return $MN$ through \texttt{\$v0} and the supplied continuation.

\section{Follow the ordinary dot product}

The ordinary entry, \texttt{matmul\_naive}, visits output row $i$, column $j$,
then reduction index $q$. Its input addresses are
\[
  A+4(iK+q),\qquad B+4(qN+j).
\]
During a dot product, the left pointer advances by four bytes and the right by
$4N$ bytes. The output address is $C+8(iN+j)$. Mixing those two element sizes can
write plausible low words while leaving half of each output untouched.

The function initializes a register accumulator to zero for each output. In this
fragment, \texttt{\$t4}/\texttt{\$t5} are the input cursors,
\texttt{\$t6} is right-column byte stride $4N$, \texttt{\$t2} is $q$, and
\texttt{\$s1} is $K$:

\noindent\begin{minipage}{\linewidth}
  \MARSCodeLabel{Accumulate one output in a register}
  \begin{minted}{mars-idea}
naive_dot:
    lw $t7, 0($t4)
    lw $t8, 0($t5)
    dmul $t7, $t7, $t8
    daddu $t3, $t3, $t7
    daddiu $t4, $t4, 4
    daddu $t5, $t5, $t6
    addiu $t2, $t2, 1
    bnec $t2, $s1, naive_dot
    nop
\end{minted}
\end{minipage}

After $q$ accepted pairs, the accumulator contains exactly the first $q$ terms
of that output's sum. When $q=K$, the function stores one complete doubleword.
A separate zero-$K$ branch skips the first input load and stores the initial zero.
The output is never used as an uninitialized accumulator.

Both worked functions reserve an 80-byte frame, save nine doubleword registers,
and restore them on return. The arithmetic uses bounded word-sized indices where
sufficient and full-width additions for pointers and sums. The supplied caller
uses aligned stack address \texttt{0x7FFF0000}, making stack placement reproducible
for this cache experiment. Grading can pass another valid aligned pointer; the
callee must restore whichever one it receives.

\section{Preserve partial sums across reduction blocks}

The blocked entry groups output rows, output columns and reduction terms into
intervals. At block origins $I,J,L$, their exclusive ends are
\[
  I_e=\min(M,I+T),\quad J_e=\min(N,J+T),\quad L_e=\min(K,L+T).
\]
It visits $I\leq i<I_e$, $J\leq j<J_e$ and $L\leq q<L_e$. Every axis needs its
own clipped end. A valid output tile can still have a shorter final reduction
interval when $K$ is not a multiple of $T$.

The worked schedule is $I,J,L,i,j,q$. It first zeros the complete output, then
visits reduction blocks in increasing order. On entry to reduction block $L$,
each output in the current tile already contains the terms with $q<L$. Loading
that partial sum, adding terms through $L_e-1$, then storing it preserves the
invariant for the next block. Replacing the load with zero would discard all
earlier reduction blocks.

For the default product and $T=2$, the first reduction block contributes
\[
  \begin{bmatrix}25&28\\73&82\end{bmatrix}.
\]
The second contributes $\begin{bmatrix}33&36\\66&72\end{bmatrix}$, giving the
complete product above. A pause before the second reduction block should show
the first matrix in output storage. That is a coherent intermediate result, not
permission to return success yet.

\begin{figure}[htbp]
  \centering
  \begin{tikzpicture}[x=7mm,y=-7mm,alt={A five by three output is partitioned into tiles of edge three. The top tile has three rows and three columns, and the lower clipped tile has two rows and three columns. Along a separate seven-term reduction axis, blocks cover indices zero through two, three through five, and six alone. The last output tile and last reduction block are clipped independently.}]
    \fill[black!14] (0,3) rectangle (3,5);
    \foreach \x in {0,...,3} {\draw[black!40] (\x,0)--(\x,5);}
    \foreach \y in {0,...,5} {\draw[black!40] (0,\y)--(3,\y);}
    \draw[thick] (0,0) rectangle (3,5) (0,3)--(3,3);
    \node[font=\small,anchor=south] at (1.5,-.15) {$M=5,\ N=3$};
    \node[font=\small,anchor=north] at (1.5,5.15) {Clipped tile: $2\times3$};
    \foreach \x in {0,...,6} {\node[font=\small] at (6+\x+.5,1.5) {\x};}
    \foreach \x in {0,...,7} {\draw[black!40] (6+\x,1)--(6+\x,2);}
    \draw[thick] (6,1) rectangle (13,2) (9,1)--(9,2) (12,1)--(12,2);
    \node[font=\small,anchor=south] at (9.5,.7) {$K=7,\ T=3$};
    \node[font=\small,anchor=north] at (9.5,2.3) {Reduction lengths: 3, 3, 1};
    \node[font=\small,align=left,anchor=west] at (6,4) {Every output receives all seven terms;\\the final reduction block adds one term.};
  \end{tikzpicture}
  \caption{Output clipping and reduction clipping are separate bounds. The final output tile can contain several cells while its final reduction interval has only one term.}
  \label{fig:revision-measurement-tiles}
\end{figure}


Blocking keeps a smaller group of input and output lines in use together, but
this implementation pays for it. It repeatedly loads and stores partial output
sums, initializes the output separately, and manages clipped block loops. The
ordinary version holds each complete dot product in a register. Correct blocking
can therefore improve line reuse while increasing executed work and memory traffic.

\section{Define the measurement interval and initial state}

The public tool \texttt{tools/measure\_revision\_matmul.py} checks products while
measuring the exact supplied source. It prepares deterministic signed input words,
a sentinel-filled output, fixed shape and method, and the same stack address.
Each cache has sixteen lines of four words, LRU replacement and 256-byte capacity.
Full, direct and two-way placements observe the same memory notices.

The \emph{cold kernel} interval begins at the selected function entry and ends
before the first instruction at \texttt{matmul\_done}. It includes frame saves,
descriptor loads, zeroing, multiplication, partial sums and frame restoration.
It excludes the caller's argument preparation and later store to
\texttt{completed}. The observer starts with empty cache lines.

The \emph{warm} observation is a second call in the same stopped execution
session, preserving the cache contents left by the first call. Its counters are
reported as differences, not cumulative totals. A third call resets the cache
before entry; it must reproduce the cold result. These terms describe the modeled
data cache. JVM compilation and session startup are separate concerns and are
outside this experiment.

A separate \emph{launch} interval starts at main and ends at the same return
boundary. It includes the caller's method read and argument/call instructions.
The native cache view initially counts this broader interval. Record the boundary
alongside its values; agreeing final products do not make differently scoped
counters comparable.

\begin{table}[htbp]
  \centering
  \small
  \begin{tabular}{llrrr}
    \toprule
    Method   & Scope                       & Instructions & Notices & Full misses \\
    \midrule
    Ordinary & cold function               & 37475        & 8744    & 5388        \\
             & second call, cache retained & 37475        & 8744    & 5383        \\
             & function after reset        & 37475        & 8744    & 5388        \\
             & launch through return       & 37523        & 8745    & 5388        \\
    Blocked  & cold function               & 54984        & 12840   & 1068        \\
             & second call, cache retained & 54984        & 12840   & 1063        \\
             & function after reset        & 54984        & 12840   & 1068        \\
             & launch through return       & 55040        & 12841   & 1068        \\
    \bottomrule
  \end{tabular}
  \caption{The same sixteen-by-sixteen product with tile edge four under different scopes and initial cache states. Counts include frame and descriptor work. The retained cache saves only five misses here; resetting reproduces the cold counts exactly. Launch includes one extra method-word read.}
  \label{tab:revision-measurement-scopes}
\end{table}


External test preparation and debugger inspection do not count as executed data
accesses. Assembly and source processing are also outside these intervals.
Instruction counts describe completed machine instructions, so a pseudo-instruction
that expands to several instructions contributes several. A source line count
cannot replace that observation.

\section{Account for traffic before interpreting misses}

The cache model counts native data notices. In this aligned code, a word load or
store contributes one notice; a doubleword load or store contributes two word
notices. Thus one \texttt{sd} is one completed instruction and two data notices.
Instruction fetches are excluded from this data-cache model.

For the ordinary function, the notice count is
\[
  A_{ordinary}=40+2MKN+2MN.
\]
The constant forty contains eighteen word notices for the nine saved doublewords,
eighteen for their restoration, and four descriptor reads. Every multiplication
loads two input words; every output store supplies two word notices.

For the blocked implementation, let $S=\lceil K/T\rceil$, with $S=0$ when $K=0$.
Its count is
\[
  A_{blocked}=40+2MKN+2MN+4MNS.
\]
The $2MN$ term is its initial zeroing. For each of the $S$ reduction blocks, every
output receives a doubleword load and store, adding $4MN$ notices. This accounting
includes the empty cases because a zero $MN$ removes all matrix terms. It is a
property of these two implementations, not a required count for every valid
solution.

For $16\times16$ inputs and $T=4$, both execute 4096 scalar multiplications.
The ordinary function has 8744 notices, while the blocked function has 12840.
The additional 4096 notices are partial-output traffic. A smaller miss count must
therefore be explained alongside more total accesses, rather than presented as
less work of every kind.

\section{Interpret a controlled comparison}

\begin{table}[htbp]
  \centering
  \small
  \setlength{\tabcolsep}{5pt}
  \begin{tabular}{llrrrrr}
    \toprule
    $(M,K,N)$    & Method         & Instructions & Notices & Full & Direct & Two-way \\
    \midrule
    $(2,3,2)$    & ordinary       & 211          & 72      & 13   & 18     & 19      \\
                 & blocked, $T=4$ & 266          & 88      & 13   & 20     & 22      \\
    $(8,8,8)$    & ordinary       & 5315         & 1192    & 188  & 275    & 324     \\
                 & blocked, $T=4$ & 7064         & 1704    & 164  & 316    & 192     \\
    $(16,16,16)$ & ordinary       & 37475        & 8744    & 5388 & 4880   & 4620    \\
                 & blocked, $T=4$ & 54984        & 12840   & 1068 & 2324   & 2544    \\
                 & blocked, $T=8$ & 43928        & 10792   & 1548 & 5300   & 5004    \\
    $(15,13,11)$ & ordinary       & 20225        & 4660    & 2521 & 2434   & 2515    \\
                 & blocked, $T=4$ & 31666        & 7300    & 1015 & 1812   & 1455    \\
    \bottomrule
  \end{tabular}
  \caption{Cold function measurements of the supplied implementations. The final three columns are misses at equal 256-byte capacity, with sixteen four-word lines and LRU replacement. Input patterns, active bases and stack address are controlled; these are reference observations, not free-writing acceptance limits.}
  \label{tab:revision-measurement-counts}
\end{table}


Table~\ref{tab:revision-measurement-counts} uses the same cold function boundaries
and cache capacities. On the sixteen-by-sixteen case, four-wide blocking reduces
full-placement misses from 5388 to 1068, while instructions rise from 37475 to
54984. Its total notices rise too. At tile edge eight, the blocked loop uses fewer
instructions than at four but incurs more full-placement misses. The smallest
matrix obtains no full-placement miss reduction at either edge.

Changing placement changes the result again. For tile edge eight on the largest
case, the blocked direct-mapped count is 5300 misses, above the ordinary version's
4880. Equal byte capacity does not give every line the same set of possible
neighbors. Fully associative LRU is also not an oracle for the best replacement
decision on every stream; use the measured counts for the stated model.

An invented cost model can expose the tradeoff. If each instruction costs one
unit and every full-placement miss adds $P$ units, four-wide blocking wins the
largest case when
\[
  54984+1068P<37475+5388P,
  \qquad P>\frac{17509}{4320}\approx4.05.
\]
Those weights are assumptions. MARS supplies no miss latency or hardware-cycle
prediction here. Actual processor instruction costs, overlap, cache hierarchy and
execution environment need their own measurements. An IDE elapsed time includes
host work as well, including observation and repainting.

Treat the table as a reproducible question about these programs. Verify arithmetic
and ownership first, change one choice, repeat the same workload and scope, and
explain any counter movement from the access schedule. Keep correctness checks
separate from exact reference measurements so an equivalent implementation can
use different instructions and traffic.

\section{Transfer the reduction to an independent interface}

Implement \texttt{matvec} in \texttt{matrix-vector.asm}, supplied in the same lab
files directory. It receives a packed $M\times K$ word matrix, a $K$-word vector,
a doubleword output array and a pointer to two words $M,K$ in
\texttt{\$a0} through \texttt{\$a3}. Dimensions are zero through sixteen, inputs
have the same signed range as above, and output row $i$ is
\[
  y[i]=\sum_{q=0}^{K-1} A[i,q]x[q].
\]
Return $M$. Preserve full saved registers, frame pointer, stack pointer, inputs,
shape and inactive output. Pointers may select interior disjoint windows and may
be above four gigabytes. Read only the selected inputs, descriptor and initialized
output; a restored private stack frame is permitted. Initialize every selected
output byte during this call before any partial-sum read.

With $M=0$, access no matrix or vector bytes and produce no output stores. With
$K=0$, write $M$ zero doublewords without input reads. There is no console I/O.
The default matrix rows $[1,2,3]$ and $[4,5,6]$, with vector $[7,8,9]$, give
$[50,122]$. Derive these independently, then decide whether to accumulate one
row at a time or use another equivalent schedule. The public checks assess the
interface and protected accesses, not a particular instruction or miss count.
